\section{Sequences and Series of Functions}

\begin{definition}[Sequence of Functions]
    \label{def:seq-of-fn}%
    We say \(\pqty{f_n}_{n \in \NN} = \pqty{f_1, f_2, \ldots}\) is a \emph{sequence of \(\CC\)-valued functions on \(X \subseteq \CC\)}, if we have \(\functiondomain{f_n}{X}{\CC}\) for each \(n \in \NN\).
\end{definition}

\begin{definition}[Pointwise Convergence]
    \label{def:pointwise-convergence}%
    If, for any \(x \in Y \subseteq X \subseteq \CC\), the numerical sequence \(\pqty{f_n \pqty{x}}_{n}\) converges, then we say the sequence \emph{pointwise converges} on \(Y\) to the function
    \[
        \function{f}{Y}{\CC}{x}{\lim_{n \to \infty} f_n \pqty{x},}
    \]
    and we say \(f_n \to f\) pointwise on \(Y\).
\end{definition}

\begin{definition}[Series of Functions]
    \label{def:series-of-functions}%
    Let \(\pqty{f_n}_{n}\) be a sequence of functions. We call the infinite sum
    \[
        \sum_{n = 1}^{\infty} f_n = f_1 + f_2 + \cdots
    \]
    as a \emph{series} of \(\CC\)-valued functions on \(X\).

    If the infinite sum converges for all \(x \in Y \subseteq X\), then we define the function
    \[
        \function{f}{Y}{\CC}{x}{\sum_{n = 1}^{\infty} f_n \pqty{x}}
    \]
    as the series.
\end{definition}

If we know the integrability, continuity, or differentiability of \(f_n\), can we say anything about \(f\) in that respect? The answer is non-trivial, since we only know before that they work for finite sums -- while for infinite sums we are dealing with two limits. In other words, we are studying whether we may swap limit/differentiation/integration/infinite summation operators. For example, as we have seen before, we can only swap limits with function operations when the function is continuous.

As an example, if we want to show that \(f\) is continuous at some \(a\), then we want to show
\[
    \lim_{x \to a} \lim_{n \to \infty} f_n \pqty{x} = \lim_{n \to \infty} \lim_{x \to a} f_n \pqty{x}.
\]

But a counterexample may be constructed. For example, reducing our perspective to sequences, we define
\[
    s_{m, n} = \frac{m}{m + n}.
\]

We have
\[
    \lim_{n \to \infty} \lim_{m \to \infty} s_{m, n} = \lim_{n \to \infty} 1 = 1,
\]
while swapping limits gives
\[
    \lim_{m \to \infty} \lim_{n \to \infty} s_{m, n} = \lim_{m \to \infty} 0 = 0.
\]

We will focus on an easier case where our functions \(f_n\) are always given in the form
\[
    f_n \pqty{x} = c_n \pqty{x - a}^n
\]
where \(a\) is the \emph{centre}, and \(c_n\) as the \emph{coefficients}.

\subsection{Basics of Power Series}

\subsubsection{Definitions and Radius of Convergence}

\begin{definition}[Power Series]
    \label{def:power-series}%
    Let \(a \in \CC\) and \(\pqty{c_n}_{n \in \NN}\) be a numerical sequence in \(\CC\). We call
    \[
        \sum_{n = 0}^{\infty} c_n \pqty{x - a}^n
    \]
    a \emph{power series} \emph{centred at \(a\)}, and with \emph{coefficients} \(\pqty{c_n}_n\).
\end{definition}

\begin{example}[Taylor Series is Power Series]
    \label{ex:taylor-series}%
    If \(f\) is smooth at \(a\), then its Taylor series at \(a\) is a power series with coefficients
    \[
        c_n = \frac{f^{\pqty{n}} \pqty{a}}{n!}.
    \]
\end{example}

Every power series trivially converges at its centre where \(x = a\).

We first study the convergence of power series, i.e., the set
\[
    Y = \set{x \in \CC}{\sum_{n = 0}^{\infty} c_n \pqty{x - a}^{n} \text{ converges at }x}.
\]

\begin{lemma}[Convergence at one Point implies Absolute Convergence in Smaller Radius]
    \label{lem:convergence-one-point-implies-convergence-smaller-radius}%
    If the power series
    \[
        \sum_{n = 0}^{\infty} c_n \pqty{x - a}^{n}
    \]
    converges for some \(x\), and we have \(\abs{y - a} < \abs{x - a}\), then the series
    \[
        \sum_{n = 0}^{\infty} c_n \pqty{y - a}^{n}
    \]
    converges absolutely.
\end{lemma}

\begin{proof}
    We conduct a comparison test to geometric series. Let
    \[
        r = \frac{\abs{y - a}}{\abs{x - a}} < 1.
    \]

    Then the power series about \(y\) satisfies
    \[
        \abs{c_n \pqty{y - a}^{n}} \leq \abs{c_n \pqty{x - a}^{n}} \cdot \abs{\frac{y - a}{x - a}}^n = \abs{c_n \pqty{x - a}^{n}} \cdot r^n
    \]
    for \(r < 1\).

    Now, since the series converges at \(x\), we have
    \[
        \lim_{n \to \infty} c_n \pqty{x - a}^{n} = 0,
    \]
    and in particular this sequence is bounded., so there is some \(M > 0\), such that
    \[
        \abs{c_n \pqty{x - a}^{n}} \leq M.
    \]

    Hence,
    \[
        \abs{c_n \pqty{y - a}^n} \leq M r^n
    \]
    for \(r < 1\), and comparing to the geometric series, we see that the series
    \[
        \sum_{n = 0}^{\infty} \abs{c_n \pqty{y - a}^{n}}
    \]
    converges, and hence the desired series is absolutely convergent.
\end{proof}

\begin{proposition}[Existence of Radius of Convergence]
    \label{prop:existence-radius-convergence}%
    Let \(\sum_{n = 0}^{\infty} c_n \pqty{x - a}^{n}\) be a power series. Then there is some \(R \in \ccintv{0}{\plusinfty}\), which is known as the \emph{radius of convergence} of the power series, such that
    \begin{itemize}
        \item the series converges absolutely for \(\abs{x - a} < R\); and
        \item the series diverges for \(\abs{x - a} > R\).
    \end{itemize}
\end{proposition}

\begin{remark}
    On the circle \(\abs{x - a} = R\) on the complex plane (or two points in \(\RR\)), we do not claim any behaviour on the convergence/divergence of those points -- it may diverge, converge absolutely, or converge conditionally on different points on the circle.
\end{remark}

If \(R = 0\), this means the series only converges at \(x = a\), and if \(R = \plusinfty\), this means the series converges absolutely for any \(x \in \CC\).

\begin{proof}
    We look at the set
    \[
        A = \set{r \geq 0}{\exists x \in \CC, \abs{x - a} = r: \text{ the power series converges at }x}.
    \]

    Trivially, \(0 \in A\). This tells us that \(A\) is non-empty. Set
    \[
        R = \sup A \in \ccintv{0}{\plusinfty}.
    \]

    If \(R = \plusinfty\), this means we have points converging at arbitrarily large radius, so by \autoref{lem:convergence-one-point-implies-convergence-smaller-radius}, this means the power series converge absolutely for any \(x \in \CC\).

    Now, if \(R\) is finite, by definition of the supremum, we know that the power series diverges for any point strictly further away from \(a\) than \(R\).
    
    For any \(\abs{y - a} < R\), then by the definition of the supremum, there is some \(r \in A\) with \(\abs{y - a} < r \leq R\). Then, some point \(x \in A\) gives convergence at radius \(r\). So since \(\abs{y - a} < r = \abs{x - a}\), by \autoref{lem:convergence-one-point-implies-convergence-smaller-radius}, the power series converges absolutely at \(y\).
\end{proof}

To compute the radius of convergence of power series, we may use the usual tests for numerical series, namely \autoref{prop:ser-test-root} (Root Test) and \autoref{prop:ser-test-ratio}.

\begin{proposition}[Root Test for Power Series]
    \label{prop:power-series-test-root}%
    Let \(\pqty{c_n}_n\) be a sequence in \(\CC\), such that
    \[
        L = \lim_{n \to \infty} \sqrt[n]{\abs{c_n}}
    \]
    exists (or is \(\plusinfty\)). Then the radius of convergence is given by
    \[
        R = \frac{1}{L}.
    \]
\end{proposition}

\begin{proposition}[Ratio Test for Power Series]
    \label{prop:power-series-test-ratio}%
    Let \(\pqty{c_n}_n\) be a sequence in \(\CC\), such that
    \[
        L = \lim_{n \to \infty} \abs{\frac{c_{n + 1}}{c_n}}
    \]
    exists (or is \(\plusinfty\)). Then the radius of convergence is given by
    \[
        R = \frac{1}{L}.
    \]
\end{proposition}

\begin{remark}
    If we have the `upgraded' version of the ratio test using \(\limsup\), then changing the \(L\) above to use \(\limsup\) will guarantee its existence, and could be used as an alternative definition.
\end{remark}

\begin{examples}[Radius of Convergence]
    \label{ex:radius-of-convergence}%
    \begin{itemize}
        \item The series
        \[
            \sum_{n = 0}^{\infty} \frac{x^n}{n!}
        \]
        converges absolutely on all of \(\CC\), i.e., \(R = \plusinfty\). Indeed, the coefficients are given by \(c_n = \frac{1}{n!}\), and \autoref{prop:power-series-test-ratio} gives
        \[
            L = \lim_{n \to \infty} \frac{\flatfrac{1}{\pqty{n + 1}!}}{\flatfrac{1}{n!}} = \lim_{n \to \infty} \frac{1}{n + 1} = 0.
        \]
        
        \item The series
        \[
            \sum_{n = 0}^{\infty} n! x^n
        \]
        has a zero radius of convergence.

        \item We investigate the convergence of the series
        \[
            \sum_{n = 1}^{\infty} \frac{x^n}{n^2}.
        \]

        The ratio test gives
        \[
            \lim_{n \to \infty} \abs{\frac{\frac{1}{\pqty{n + 1}^2}}{\frac{1}{n^2}}} = 1
        \]
        which gives \(R = 1\). Therefore, the power series converges absolutely with \(\abs{x} < 1\), and diverges with \(\abs{x} > 1\).

        For \(\abs{x} = 1\), we note that
        \[
            \sum_{n = 1}^{\infty} \abs{\frac{x^n}{n^2}} = \sum_{n = 1}^{\infty} \frac{1}{n^2} = \frac{\pi^2}{6}
        \]
        converges absolutely.

        Therefore, this power series converges absolutely for \(\abs{x} \leq 1\), and diverges \(\abs{x} > 1\).

        \item Now investigating the series
        \[
            \sum_{n = 1}^{\infty} \frac{x^n}{n}.
        \]

        The radius of convergence remains to be \(1\). However, with \(\abs{x} = 1\), the behaviour is more subtle.

        If \(x = 1\), this gives divergence since this is the harmonic series. However, if \(x = -1\), this gives (conditional) convergence since this is the alternating harmonic series.
    \end{itemize}
\end{examples}

\begin{remark}
    The set of convergence points
    \[
        Y = \set{x \in \CC}{\sum_{n = 0}^{\infty} c_n \pqty{x - a}^n \text{ converges}}
    \]
    and hence gives a possible domain for the function of the power series. It must contain the set
    \[
        \set{x \in \CC}{\abs{x - a} < R}.
    \]

    However, it may or may not contain the points
    \[
        \set{x \in \CC}{\abs{x - a} = R},
    \]
    and definitely does not contain
    \[
        \set{x \in \CC}{\abs{x - a} > R}.
    \]
\end{remark}

\subsubsection{Analytical Properties of Power Series}

Given that the terms of the sums are polynomials, this means they are continuous, integrable, and differentiable. We would like to study these properties of the infinite sum as well. Formally, if we have the function
\[
    \function{f}{Y}{\CC}{x}{\sum_{n = 0}^{\infty} c_n \pqty{x - a}^n,}
\]
then for example, can we conclude that
\[
    \dv{x} f\pqty{x} = \sum_{n = 0}^{\infty} c_n \dv{x} \pqty{x - a}^n,
\]
and can we conclude that
\[
    \int f\pqty{x} \dd{x} = \sum_{n = 0}^{\infty} c_n \int \pqty{x - a}^{n} \dd{x}?
\]

\begin{example}[Differentiating Power Series Term-by-Term]
    \label{ex:differentiating-power-series-term-by-term}%
    \begin{itemize}
        \item The power series \(\sum_{n = 0}^{\infty} \frac{x^n}{n!}\) converges on all of \(\CC\). Then we have
        \[
            \sum_{n = 0}^{\infty} \frac{1}{n!} \dv{x} \pqty{x^n} = \sum_{n = 1}^{\infty} \frac{x^{n - 1}}{\pqty{n - 1}!} = \sum_{n = 0}^{\infty} \frac{x^n}{n!}
        \]
        convergences on \(\CC\).

        \item The power series \(\sum_{n = 1}^{\infty} \frac{x^n}{n^2}\) converges on the unit closed disk. Differentiating term-by-term, we have
        \[
            \sum_{n = 1}^{\infty} \frac{1}{n^2} \dv{x} \pqty{x^n} = \sum_{n = 1}^{\infty} \frac{x^{n - 1}}{n}.
        \]

        Although this has the same radius of convergence (which is \(1\)), it does not converge on the whole of the \emph{closed} unit disk -- but it does converge on the open unit disk.
    \end{itemize}
\end{example}

Clearly, we need to be careful: term-by-term operations might not hold on the \emph{entire} set of convergence of power series. But we may show that they do converge within the radius of convergence. We first introduce two notations to help us. 

\begin{definition}[Open Ball and Closed Disk]
    \label{def:open-ball-closed-disk}%
    Let \(a \in \CC\) be the \emph{centre} and \(R > 0\) be the \emph{radius}. Then the \emph{open ball} with centre \(a \in \CC\) and radius \(R > 0\) is defined as
    \[
        B_R \pqty{a} = \set{x \in \CC}{\abs{x - a} < R}
    \]
    and the \emph{closed disk} as
    \[
        D_R \pqty{a} = \set{x \in \CC}{\abs{x - a} \leq R}.
    \]
\end{definition}

\begin{proposition}[Continuity of Power Series]
    \label{prop:continuity-of-power-series}%
    Let
    \[
        \sum_{n = 0}^{\infty} c_n \pqty{x - a}^{n}
    \]
    be a power series with a non-zero radius of convergence \(R > 0\). Then the function
    \[
        \function{f}{B_R \pqty{a}}{\CC}{x}{\sum_{n = 0}^{\infty} c_n \pqty{x - a}^n}
    \]
    is continuous inside \(D_r\pqty{a}\) for every \(r < R\). In particular, it is continuous on the whole of the open unit ball.
\end{proposition}

\begin{proof}[\proofname\ (Non-Examinable)]
    We introduce the notation
    \[
        S_N\pqty{x} = \sum_{n \leq N} c_n \pqty{x - a}^n, \quad T_N\pqty{x} = \sum_{n > N} c_n \pqty{x - a}^n
    \]
    and notice the tail may be bounded above by
    \[
        \abs{T_N\pqty{a}} \leq \sum_{n > N} \abs{c_n} \abs{x - a}^n \leq \sum_{n > N} \abs{c_n} r^n
    \]
    for all \(x \in D_r\pqty{a}\) with \(r < R\). By the definition of radius of convergence, for all \(\epsilon > 0\), there is some \(N = N \pqty{\epsilon} \in \NN\), such that
    \[
        \sum_{n > N} \abs{c_n} r^n < \epsilon,
    \]
    and hence
    \[
        \forall x \in D_r\pqty{a}: \abs{T_N\pqty{x}} < \epsilon
    \]

    Take \(x_0 \in D_r\pqty{a}\). We want to show that \(f\) is continuous at \(x_0\). By the triangular inequality, we have
    \begin{align*}
        \abs{f\pqty{x} - f\pqty{x_0}} &\leq \abs{S_N\pqty{x} - S_N\pqty{x_0}} + \abs{T_N\pqty{x}} + \abs{T_N \pqty{x_0}} \\
        &\leq \abs{S_N\pqty{x} - S_N\pqty{x_0}} + 2 \epsilon
    \end{align*}
    for all \(x \in D_r\pqty{a}\). To conclude the proof, we use the fact that polynomials are continuous, and choose some \(\delta = \delta\pqty{\epsilon}\) such that
    \[
        \abs{S_N\pqty{x} - S_N\pqty{x_0}} < \epsilon.
    \]

    This finishes the proof.
\end{proof}

\begin{proposition}[Integrability of Real Power Series, Term-by-Term]
    \label{prop:integrability-of-power-series-term-by-term}%
    Let
    \[
        \sum_{n = 0}^{\infty} c_n \pqty{x - a}^n
    \]
    be a real power series, with radius of convergence \(R > 0\). Then the function
    \[
        \function{f}{\oointv{a - R}{a + R}}{\RR}{x}{\sum_{n = 0}^{\infty} c_n \pqty{x - a}^n}  
    \]
    is integrable on every \(\ccintv{a - r}{a + r}\) for every \(r < R\), and
    \[
        \int_{a}^{x} f\pqty{t} \dd{t} = \sum_{n = 0}^{\infty} c_n \frac{\pqty{x - a}^{n + 1}}{n + 1}
    \]
    for any \(a - R < x < a + R\).
\end{proposition}

\begin{proof}[\proofname\ (Non-Examinable)]
    Continuity implies integrability, so it remains to show the formula for integration. We use the same idea, by splitting up the power series into a polynomial and a tail. We have
    \begin{align*}
        \abs{\int_{a}^{x} f\pqty{t} \dd{t} - \sum_{n \leq N} \frac{c_n \pqty{x - a}^{n + 1}}{n + 1}} &= \abs{\int_{a}^{x} f\pqty{t} \dd{t} - \int_{a}^{x} S_N\pqty{t} \dd{t}} \\
        &= \abs{\int_{a}^{x} T_N\pqty{t} \dd{t}} \\
        &\leq \sup_{\abs{t - a} < r} \abs{T_N \pqty{t}} \abs{\int_{a}^{x} \dd{t}} \\
        &\leq r \sup_{\abs{t - a} < r} \abs{T_N \pqty{t}} \\
        &\leq r \epsilon
    \end{align*}
    for some sufficiently large \(N = N\pqty{\epsilon}\). Hence,
    \[
        \int_{a}^{x} f\pqty{t} \dd{t} = \lim_{n \to \infty} \sum_{n = 0}^{N} \frac{c_n \pqty{x - a}^{n + 1}}{n + 1} = \sum_{n = 0}^{\infty} \frac{c_n \pqty{x - a}^{n + 1}}{n + 1}.
    \]
\end{proof}

\begin{proposition}[Differentiability of Power series, Term-by-Term]
    \label{prop:differentiability-of-power-series-term-by-term}%
    Let
    \[
        \sum_{n = 0}^{\infty} c_n \pqty{x - a}^{n}
    \]
    be a power series, with a non-zero radius of convergence \(R > 0\). Then the function
    \[
        \function{f}{B_R \pqty{a}}{\CC}{x}{\sum_{n = 0}^{\infty} c_n \pqty{x - a}^n}
    \]
    is differentiable on the whole of its domain, and we have
    \[
        f'\pqty{x} = \sum_{n = 1}^{\infty} n c_n \pqty{x - a}^{n - 1}.
    \]
\end{proposition}

\begin{proof}[\proofname\ (Non-Examinable, Using Integrability, Real Power Series)]
    We let \(g\) be the series we claim to be the derivative. If \(g\) has radius of convergence at least \(R\), then \(g\) is continuous on \(\ccintv{a - r}{a + r}\) for all \(r < R\) by \autoref{prop:continuity-of-power-series} (Continuity of Power Series), and using \autoref{prop:integrability-of-power-series-term-by-term} (Integrability of Power Series), we have
    \[
        \int_{a}^{x} g\pqty{t} \dd{t} = \sum_{n = 1}^{\infty} n c_n \frac{\pqty{x - a}^{n}}{n} = f\pqty{x} - c_0.
    \]

    Hence, \autoref{thm:ftc-1} (Fundamental Theorem of Calculus), we have
    \[
        f'\pqty{x} = g\pqty{x}
    \]
    for all \(x \in \ccintv{a - r}{a + r}\) for all \(r < R\), and hence
    \[
        f'\pqty{x} = g\pqty{x}
    \]
    for all \(x \in \oointv{a - R}{a + R}\).

    Now, it remains to show that the new series has radius of convergence at least \(R\). Take \(r < R\), and let \(s \in \oointv{r}{R}\). We have
    \begin{align*}
        \abs{n c_n \pqty{x - a}^{n - 1}} &= \abs{c_n} s^n \frac{n}{s} \pqty{\frac{\abs{x - a}}{s}}^{n - 1} \\
        &\leq \abs{c_n} s^n \frac{n}{s} \pqty{\frac{r}{s}}^{n - 1}
    \end{align*}
    for all \(\abs{x - a} < r\). Since \(s < R\), \(\abs{c_n} s^n \to \infty\), and in particular, \(\abs{c_n} s^n\) is bounded. Hence,
    \[
        \sum_{n = 0}^{\infty} M \frac{n}{s} \pqty{\frac{r}{s}}^{n - 1}
    \]
    converges based on the assumption \(\frac{r}{s} < 1\). Hence, the series
    \[
        \sum_{n = 1}^{\infty} n c_n \pqty{x - a}^{n - 1}
    \]
    converges also in \(\abs{x - a} < r\), for all \(r < R\).
\end{proof}

\begin{proof}[\proofname\ (Non-Examinable, Direct, Complex Power Series)]
    Consider \(x_0 \in B_R\pqty{a}\). Let \(\abs{x_0 - a} < r < R\) and \(\abs{x - a} < r\).

    Recall the identity
    \[
        \frac{\pqty{x - a}^n - \pqty{x_0 - a}^n}{x - x_0} = \sum_{k = 0}^{n - 1} \pqty{x - a}^k \pqty{x_0 - a}^{n - 1 - k}.
    \]

    Subtracting the target derivative term \(n \pqty{x - a}^{n - 1}\), we have
    \begin{align*}
        \frac{\pqty{x - a}^n - \pqty{x_0 - a}^n}{x - x_0} - n \pqty{x_0 - a}^{n - 1} &= \sum_{k = 0}^{n - 1} \pqty{\pqty{x - a}^k \pqty{x_0 - a}^{n - 1 - k} - \pqty{x_0 - a}^{n - 1}} \\
        &= \sum_{k = 0}^{n - 1} \pqty{x_0 - a}^{n - 1 - k} \pqty{\pqty{x - a}^k - \pqty{x_0 - a}^k} \\
        &= \pqty{x - x_0} \sum_{k = 0}^{n - 1} \pqty{x_0 - a}^{n - 1 - k} \sum_{j = 0}^{k - 1} \pqty{x - a}^{j} \pqty{x_0 - a}^{k - 1 - j}.
    \end{align*}

    We have
    \[
        \abs{x - a}, \abs{x_0 - a} < r,
    \]
    and hence
    \[
        \abs{\frac{\pqty{x - a}^n - \pqty{x_0 - a}^n}{x - x_0} - n \pqty{x - a}^{n - 1}} \leq \abs{x - x_0} \binom{n}{2} r^{n - 2}.
    \]

    We may check that
    \[
        \sum_{n = 2}^{\infty} \abs{c_n} \binom{n}{2} r^{n - 2}
    \]
    converges, and let this sum be \(M\).
    
    Now we investigate the difference of the power series. Since both series converge, we may subtract them term-by-term, giving
    \[
        \frac{f\pqty{x} - f\pqty{x_0}}{x - x_0} - \sum_{n = 1}^{\infty} n c_n \pqty{x_0 - a}^{n - 1} = \sum_{n = 2}^{\infty} c_n \pqty{\frac{\pqty{x - a}^n - \pqty{x_0 - a}^n}{x - x_0} - n \pqty{x_0 - a}^{n - 1}},
    \]
    and hence
    \[
        \abs{\frac{f\pqty{x} - f\pqty{x_0}}{x - x_0} - \sum_{n = 1}^{\infty} n c_n \pqty{x_0 - a}^{n - 1}} \leq \abs{x - x_0} \sum_{n = 2}^{\infty} \abs{c_n} \binom{n}{2} r^{n - 2} = M \abs{x - x_0}.
    \]

    Taking \(x \to x_0\) gives us the desired result.
\end{proof}

\begin{examples}[Differentiating and Integrating Power Series]
    \label{ex:differentiating-integrating-power-series}%
    \begin{itemize}
        \item Consider \(\sum_{n = 0}^{\infty} \frac{x^n}{n!}\) which converges on all of \(\CC\). Term-by-term differentiation and integration gives us
        \[
            \dv{x} \sum_{n = 0}^{\infty} \frac{x^n}{n!} = \sum_{n = 0}^{\infty} \dv{x} \pqty{x^n} \frac{1}{n!} = \sum_{n = 1}^{\infty} \frac{x^{n - 1}}{\pqty{n - 1}!} = \sum_{n = 0}^{\infty} \frac{x^n}{n!}
        \]
        and
        \[
            \int_{0}^{x} \sum_{n = 0}^{\infty} \frac{t^n}{n!} \dd{t} = \sum_{n = 0}^{x} \int_{0}^{x} \frac{t^n}{n!} \dd{t} = \sum_{n = 0}^{\infty} \frac{x^{n + 1}}{\pqty{n + 1}!} = \sum_{n = 0}^{\infty} \frac{x^n}{n!} - 1.
        \]
        \item Consider \(\sum_{n = 1}^{\infty} \frac{x^n}{n^2}\) has radius of convergence \(1\). For all \(\abs{x} < 1\), we have
        \[
            \dv{x} \sum_{n = 1}^{\infty} \frac{x^n}{n^2} = \sum_{n = 1}^{\infty} \frac{x^{n - 1}}{n}
        \]
        and
        \[
            \int_{0}^{x} \sum_{n = 1}^{\infty} \frac{t^n}{n^2} \dd{t} = \sum_{n = 1}^{\infty} \frac{x^{n + 1}}{n^2 \pqty{n + 1}}.
        \]
    \end{itemize}
\end{examples}

\subsection{Exponentials and Logarithms}

\begin{definition}[Exponential Function]
    \label{def:exponential-function}%
    We define the \emph{exponential function}, denoted as \(\exp\), as
    \[
        \function{\exp}{\CC}{\CC}{z}{\sum_{n = 0}^{\infty} \frac{z^n}{n!}.}
    \]
\end{definition}

\begin{remarks}
    This definition of the exponential agrees with our high-school understanding of \(e^x\) in the following ways.
    \begin{itemize}
        \item Since \(e^x\) differentiates to \(e^x\), evaluating the derivatives at \(x = 0\) tells us that the Taylor series would be equal to
        \[
            \sum_{n = 0}^{\infty} \frac{z^n}{n!}
        \]
        which agrees with our definition of \(\exp\).

        \item Using \autoref{thm:taylor-thm-langrange-remainder} (Taylor's Theorem) with Lagrange's Remainder, we may evaluate the remainder term as
        \[
            R_{n, 0} \pqty{x} = \frac{e^c x^n}{n!}
        \]
        for some \(c \in \ccintv{0}{x}\). Hence,
        \[
            \abs{R_{n, 0} \pqty{x}} = \abs{e^c} \frac{\abs{x}^n}{n!} \to 0
        \]
        as \(n \to \infty\). Hence,
        \[
            e^x = \sum_{n = 0}^{\infty} \frac{x^n}{n!}
        \]
        for all real number \(x \in \RR\).
    \end{itemize}
\end{remarks}

\begin{lemma}[Properties of Exponential Function]
    \label{lem:properties-exp}%
    We have
    \begin{itemize}
        \item \(\exp\) is smooth, with \(\exp'\pqty{z} = \exp\pqty{z}\);
        \item \(\exp\pqty{0} = 1\); and
        \item \(\exp\pqty{a + b} = \exp\pqty{a} \cdot \exp\pqty{b}\).
    \end{itemize}
\end{lemma}

\begin{proof}
    The first two properties arise from differentiating/dividing the power series.

    For the third property, we set
    \[
        f_{a, b} \pqty{z} = \exp\pqty{a + b - z} \cdot \exp\pqty{z}.
    \]

    Then we have
    \begin{align*}
        f'\pqty{z} &= - \exp'\pqty{a + b - z} \exp\pqty{z} + \exp\pqty{a + b - z} \exp'\pqty{z} \\
        &= - f\pqty{z} + f\pqty{z} \\
        &= 0
    \end{align*}
    and hence \(f\) is constant on \(\CC\). Therefore,
    \[
        \exp\pqty{a} \exp\pqty{b} = f\pqty{b} = f\pqty{0} = \exp\pqty{a + b}.
    \]
\end{proof}

\begin{lemma}[Properties of Exponential Function on \(\RR\)]
    \label{lem:properties-exp-r}%
    Consider the restriction of \(\exp\) on \(\RR\)
    \[
        \functiondomain{\rightbar{\exp}_{\RR}}{\RR}{\RR}.
    \]

    Then
    \begin{itemize}
        \item \(\exp\) is smooth, with \(\exp'\pqty{x} = \exp\pqty{x}\);
        \item \(\exp\pqty{x + y} = \exp\pqty{x} \cdot \exp\pqty{y}\);
        \item \(\exp\pqty{x} > 0\) for all \(x \in \RR\);
        \item \(\exp \pqty{x} \to \plusinfty\) as \(x \to \plusinfty\), and \(\exp \pqty{x} \to \minusinfty\) as \(x \to \minusinfty\);
        \item \(\exp\) is strictly increasing;
        \item \(\functiondomain{\exp}{\RR}{\oointv{0}{\plusinfty}}\) is a bijection.
    \end{itemize}
\end{lemma}

\begin{proof}
    The first two properties follow from \autoref{lem:properties-exp}.

    For \(x > 0\), note that every term in the infinite series is positive, and in particular,
    \[
        \exp\pqty{x} \geq 1 + x > 0
    \]
    for all \(x > 0\), with \(\exp\pqty{x} \to \plusinfty\) as \(x \to \plusinfty\).

    On the other hand, for \(x < 0\), note we have
    \[
        \exp\pqty{-x} \cdot \exp\pqty{x} = \exp\pqty{0} = 1,
    \]
    so since \(\exp\pqty{-x} > 0\), we must have \(\exp\pqty{x} > 0\). Furthermore,
    \[
        \exp\pqty{x} = \frac{1}{\exp\pqty{-x}} \to 0
    \]
    as \(x \to \minusinfty\), since \(\exp\pqty{-x} \to \plusinfty\).

    Since \(\exp'\pqty{x} = \exp\pqty{x} > 0\), it is strictly increasing, and hence injective into \(\oointv{0}{\plusinfty}\). Now given any
    \[
        y \in \oointv{0}{\plusinfty} = \oointv{\lim_{x \to \minusinfty} \exp\pqty{x}}{\lim_{x \to \plusinfty} \exp\pqty{x}},
    \]
    we have some \(a, b \in \RR\), with \(\exp a < y < \exp b\), so by \autoref{thm:intermediate-value} (Intermediate Value Theorem), there is some \(x \in \ccintv{a}{b}\) such that \(\exp\pqty{x} = y\). This shows that it is surjective, hence finishing the proof.
\end{proof}

Since \(\exp\) is a bijection, it must have some inverse function.

\begin{definition}[Logarithm Function]
    \label{def:log}%
    The \emph{logarithm function}, denoted as \(\log\), is defined as the inverse of the exponential function.

    \[
        \functiondomain{\log}{\oointv{0}{\plusinfty}}{\RR}.
    \]
\end{definition}

\begin{lemma}[Properties of Logarithm Functions]
    \label{lem:properties-log}%
    We have
    \begin{itemize}
        \item \(\log\) is bijective with \(\log \exp x = x\) for all \(x \in \RR\), and \(\exp \log x = x\) for all \(x \in \oointv{0}{\plusinfty}\);
        \item \(\log\pqty{yz} = \log y + \log z\);
        \item \(\log\) is smooth and strictly increasing, with \(\log' \pqty{x} = \frac{1}{x}\);
        \item \(\log\pqty{1} = 0\), and \(\log y = \int_1^x \frac{1}{t} \dd{t}\);
        \item \(\log y \to \plusinfty\) as \(y \to \plusinfty\), and \(\log y \to \minusinfty\) as \(y \to 0^+\).
    \end{itemize}
\end{lemma}

\begin{proof}
    The first property is by definition.

    For any \(y = \exp u \in \oointv{0}{\plusinfty}\) and \(z = \exp v \in \oointv{0}{\plusinfty}\), we have
    \[
        \log \pqty{yz} = \log \pqty{\exp u \cdot \exp v} = \log \pqty{\exp \pqty{u + v}} = u + v = \log y + \log z
    \]
    as desired.

    For the third property, we may apply \autoref{thm:inverse-fn-thm-differentiable} (Inverse Function Theorem) to get
    \[
        \log'\pqty{y} = \frac{1}{\exp'\pqty{\log y}} = \frac{1}{y}.
    \]

    For the fourth property, the first point comes from \(\exp 0 = 1\), and hence by \autoref{thm:ftc-1} (Fundamental Theorem of Calculus), we have
    \[
        \log y = \int_1^y \log' t \dd{t} = \int_1^y \frac{1}{t} \dd{t}.
    \]

    The final part inherits from the limit properties of the exponential function.
\end{proof}

\begin{remark}
    We have, for \(y \in \oointv{-1}{1}\), that
    \begin{align*}
        \log\pqty{1 + y} &= \int_0^y \frac{\dd{t}}{1 + t} \\
        &= \int_0^y \frac{\dd{t}}{1 - \pqty{-t}} \\
        &= \int_0^y \sum_{n = 0}^{\infty} \pqty{-1}^n t^n \dd{t} \\
        &= \sum_{n = 0}^{\infty} \int_0^y \pqty{-1}^n t^n \dd{t} \\
        &= \sum_{n = 0}^{\infty} \frac{\pqty{-1}^{n} y^{n + 1}}{n + 1}.
    \end{align*}

    In particular, using Analysis II techniques (Abel's Theorem), or quite some tedious integral comparison, we may show that in the case for \(y = 1\), that
    \[
        \log 2 = \sum_{n = 1}^{\infty} \frac{\pqty{-1}^{n - 1}}{n}.
    \]

    We have
    \begin{align*}
        \abs{\log 2 - \sum_{n = 0}^{N} \int_{0}^{1} \pqty{-1}^n t^n \dd{t}} &= \abs{\int_{0}^{§} \frac{\dd{t}}{1 + t} - \int_{0}^{1} \frac{1 - \pqty{-t}^{N + 1}}{1 + t} \dd{t}} \\
        &\leq \int_{0}^{1} \abs{\frac{1}{1 + t} - \frac{1 - \pqty{-t}^{N + 1}}{1 + t}} \dd{t} \\
        &\leq \int_{0}^{1} t^{N + 1} \dd{t} \\
        &= \frac{1}{N + 2} \\
        &\to 0
    \end{align*}
    as \(N \to \infty\)
\end{remark}

Now, we are ready to define more general powers.

\begin{definition}[Powers]
    \label{def:general-powers}%
    We define for \(x > 0\), the \emph{base}, and \(\alpha \in \RR\), the \emph{exponent}, that
    \[
        \Gamma_{\alpha} \pqty{x} = \exp \pqty{\alpha \log x}.
    \]
\end{definition}

\begin{remark}
    We will soon see that this is consistent with \(x^\alpha\) defined before using elementary methods, in particular, for \(\alpha \in \QQ\), and for \(\exp x\) and \(e^x\) for \(x \in \RR\).
\end{remark}

\begin{lemma}[Properties of Powers]
    \label{lem:properties-power}%
    For any \(x, y > 0, \alpha, \beta \in \RR\), we have
    \begin{itemize}
        \item \(\Gamma_\alpha \pqty{xy} = \Gamma_\alpha \pqty{x} \Gamma_\alpha \pqty{y}\);
        \item \(\Gamma_{\alpha + \beta} \pqty{x} = \Gamma_\alpha \pqty{x} \Gamma_\beta \pqty{x}\);
        \item \(\Gamma_\alpha \pqty{\Gamma_\beta \pqty{x}} = \Gamma_{\alpha \beta} \pqty{x}\);
        \item \(\Gamma_1\pqty{x} = x\) and \(\Gamma_0 \pqty{x} = 1\).
    \end{itemize}
\end{lemma}

\begin{proof}
    The first, second and fourth property follows directly from the properties of the exponential and the logarithm.

    For the third property, we have
    \begin{align*}
        \Gamma_\alpha \pqty{\Gamma_\beta \pqty{x}} &= \exp\pqty{\alpha \log\pqty{\exp \pqty{\beta \log x}}} \\
        &= \exp\pqty{\alpha \beta \log x} \\
        &= \Gamma_{\alpha \beta} \pqty{x}
    \end{align*}
    as desired.
\end{proof}

As a corollary, we see that for natural numbers \(p, q \in \NN\) and \(x > 0\), we have
\[
    \Gamma_p \pqty{x} = \underbrace{\Gamma_1 \pqty{x} \cdots \Gamma_1 \pqty{x}}_{p \text{ terms}} = x^p,
\]
and
\[
    \Gamma_{-p} \pqty{x} = \underbrace{\Gamma_{-1} \pqty{x} \cdots \Gamma_{-1} \pqty{x}}_{p \text{ terms}} = x^{-p}.
\]

Furthermore, we have
\[
    \pqty{\Gamma_{\flatfrac{1}{p}} \pqty{x}}^{p} = \underbrace{\Gamma_{\flatfrac{1}{p}} \pqty{x} \cdots \Gamma_{\flatfrac{1}{p}} \pqty{x}}_{p \text{ terms}} = \Gamma_{\flatfrac{1}{p} + \cdots + \flatfrac{1}{p}} \pqty{x} = \Gamma_1 \pqty{x} = x,
\]
and hence
\[
    \Gamma_{\flatfrac{1}{p}} \pqty{x} = x^{\flatfrac{1}{p}}.
\]

Hence,
\[
    x^{\flatfrac{p}{q}} = \Gamma_{\flatfrac{p}{q}} \pqty{x},
\]
and we have
\[
    \Gamma_{\alpha} \pqty{x} = x^\alpha
\]
for all \(\alpha \in \QQ\).

As for the agreement with the exponential function, we have
\[
    \exp\pqty{x} = \exp\pqty{x \log e} = \Gamma_x \pqty{e} = e^x.
\]

We now study some limit properties on monomials, exponentials and logarithms.

\begin{proposition}[Limit Properties of Monomials, Exponentials and Logarithms]
    \label{prop:lim-prop-monomial-exponential-logarithm}%
    We have the following properties.
    \begin{enumerate}
        \item \emph{Exponentials dominate Monomials.} For all \(r \in \RR\), we have \(x^r \exp\pqty{-x} \to 0\) as \(x \to \plusinfty\).
        \item \emph{Monomials dominate Logarithms.} For all \(r > 0\), we have \(x^{-r} \log x \to 0\) as \(x \to \plusinfty\), and \(x^r \log x \to 0\) as \(x \to 0^+\).
    \end{enumerate}
\end{proposition}

\begin{proof}
    For the first limit, we let \(x > 1\). Then, since \(x > 0\), we have
    \[
        \exp x = \sum_{k = 0}^{\infty} \frac{x^k}{k!} > \frac{x^n}{n!},
    \]  
    for all \(n \in \NN\).

    We pick \(n\) such that \(n - r \geq 1\), then
    \[
        0 \leq \frac{x^r}{\exp x} \leq \frac{n!}{x^{n - r}} \leq \frac{n!}{x}.
    \]

    Taking the limit as \(x \to \plusinfty\), and using the sandwich theorem gives us the result as desired.

    For the second limit, we have \(t^{-1} \leq t^{\epsilon - 1}\) for all \(t \geq 1\) and \(\epsilon > 0\). Pick \(\epsilon\) such that \(r > \epsilon\), then
    \[
        0 \leq x^{-r} \log x = x^{-r} \int_{1}^{x} \frac{\dd{t}}{t} \leq x^{-r} \int_{1}^{x} t^{\epsilon - 1} \dd{t} = \frac{x^{\epsilon - r}}{\epsilon}
    \]
    and take \(x \to \plusinfty\) with the sandwich theorem gives the result as desired.

    For the limit where \(x \to 0^+\), we note
    \begin{align*}
        \lim_{x \to 0^+} \pqty{x^r \log x} &= - \lim_{t \to \plusinfty} \pqty{t \exp \pqty{- rt}} && x = e^{-t} \\
        &= - \lim_{y \to \plusinfty} \pqty{\frac{y e^{-y}}{r}} && y = rt \\
        &= 0
    \end{align*}
    by the first limit.
\end{proof}

\begin{proposition}[Important Limit involving \(e\)]
    We have
    \[
        \pqty{1 + \frac{a}{n}}^n \to \exp \pqty{a}
    \]
    as \(n \to \infty\), for all \(a \in \RR\).
\end{proposition}

\begin{proof}
    We consider \(\log\pqty{1 + x}\) on \(\ccintv{0}{\flatfrac{a}{n}}\), where \(n > \abs{a}\). By \autoref{thm:mean-value-thm} (Mean Value Theorem), there is some \(c \in \oointv{0}{\flatfrac{a}{n}}\) (or the other way round for the interval if \(a < 0\)), such that
    \[
        \rightbar{\dv{x}}_{x = c} \log \pqty{1 + x} = \frac{\log\pqty{1 + \flatfrac{a}{n}} - \log\pqty{1 + 0}}{\flatfrac{a}{n} - 0},
    \]
    which means
    \[
        \frac{1}{1 + c} = \frac{\log \pqty{1 + \flatfrac{a}{n}}}{\flatfrac{a}{n}},
    \]
    so
    \[
        n \log \pqty{1 + \frac{a}{n}} = \frac{a}{1 + c}
    \]
    for some \(c \in \ccintv{0}{\flatfrac{a}{n}}\).

    We want to show that
    \[
        \pqty{1 + \frac{a}{n}}^n \to \exp\pqty{a},
    \]
    and it is equivalent to show that
    \[
        n \log\pqty{1 + \frac{a}{n}} \to a.
    \]

    As \(n \to \infty, c \in \oointv{0}{\flatfrac{a}{n}}\) so \(c \to 0\), and hence
    \[
        n \log\pqty{1 + \frac{a}{n}} = \frac{a}{1 + c} \to a
    \]
    and this shows the result as desired.
\end{proof}

\subsection{Trigonometric Functions}

\begin{definition}[Trigonometric Functions]
    \label{def:trigonometric-functions}%
    We define the \emph{cosine function} as
    \[
        \function{\cos}{\CC}{\CC}{z}{\frac{1}{2} \bqty{\exp \pqty{iz} + \exp \pqty{- iz}} = \sum_{k = 0}^{\infty} \frac{\pqty{-1}^k z^{2k}}{\pqty{2k}!},}
    \]
    and the \emph{sine function} as
    \[
        \function{\sin}{\CC}{\CC}{z}{\frac{1}{2i} \bqty{\exp \pqty{iz} - \exp \pqty{- iz}} = \sum_{k = 0}^{\infty} \frac{\pqty{-1}^k z^{2k + 1}}{\pqty{2k + 1}!}}
    \]
\end{definition}

\begin{lemma}[Properties of the Trigonometric Functions]
    \label{lem:properties-trigonometric-functions}%
    We have for \(w, z \in \CC\), that:
    \begin{itemize}
        \item \(\cos 0 = 1, \sin 0 = 0\);
        \item \(\cos' z = - \sin z\), \(\sin' z = \cos z\);
        \item \(\sin \pqty{z + w} = \sin z \cos w + \cos z \sin w\);
        \item \(\cos \pqty{z + w} = \cos z \cos w - \sin z \sin w\);
        \item \(\sin^2 z + \cos^2 z = 1\).
    \end{itemize}
\end{lemma}

\begin{proof}
    These are all direct consequences of the definition using the complex exponentials.
\end{proof}

\begin{remark}
    If we restrict the functions to \(\RR\), then in particular, the final line tells us that for \(x \in \RR\)
    \[
        \abs{\sin x}, \abs{\cos x} \leq 1,
    \]
    but this is not necessarily true for all of \(\CC\).
\end{remark}

\begin{proposition}[Periodicity of Trigonometric Functions]
    \label{prop:periodicity-trigonometric-functions}%
    There is some smallest positive real number \(\pi\), such that
    \[
        \cos \pqty{\frac{\pi}{2}} = 0, \quad \sin \pqty{\frac{\pi}{2}} = 1,
    \]
    and that for all \(x \in \RR\), we have
    \begin{itemize}
        \item \(\functiondomain{\sin}{\RR}{\RR}\) and \(\functiondomain{\cos}{\RR}{\RR}\) are periodic with period \(2 \pi\);
        \item \(\sin \pqty{x + \pi} = - \sin x\), and \(\cos \pqty{x + \pi} = - \cos x\);
        \item \(\sin \pqty{x + \frac{\pi}{2}} = \cos x\), and \(\cos \pqty{x + \frac{\pi}{2}} = - \sin x\).
    \end{itemize}
\end{proposition}

\begin{proof}
    Once we showed the existence of the smallest \(\pi > 0\), the remaining three propositions follow from the additions formulae.

    We look at \(x \in \oointv{0}{2}\), since by differentiating, we have
    \[
        \cos' x = - \sin x = - \bqty{\pqty{x - \frac{x^3}{3!}} + \pqty{\frac{x^5}{5!} - \frac{x^7}{7!}} + \cdots} < 0,
    \]
    since all the terms in the brackets are positive, using that the series is absolutely convergent, and hence bracketing of terms is allowed. Thus, \(\cos\) has at most one root \(\oointv{0}{2}\). To see that the existence of the root, we have
    \[
        \cos \pqty{\sqrt{2}} = \pqty{1 - \frac{\pqty{\sqrt{2}}^2}{2!}} + \pqty{\frac{\pqty{\sqrt{2}}^4}{4!} - \frac{\pqty{\sqrt{2}}^6}{6!}} + \cdots > 0
    \]
    and
    \[
        \cos \pqty{\sqrt{3}} = \pqty{1 - \frac{\pqty{\sqrt{3}}^2}{2!} + \frac{\pqty{\sqrt{3}}^4}{4!}} - \pqty{\frac{\pqty{\sqrt{3}}^6}{6!} - \frac{\pqty{\sqrt{3}}^8}{8!}} - \cdots < 0.
    \]
    
    By intermediate value theorem, \(\cos\) has a root \(\flatfrac{\pi}{2} \in \oointv{\sqrt{2}}{\sqrt{3}} \subset \oointv{0}{2}\). Now, we have
    \[
        \sin^2\pqty{\frac{\pi}{2}} = 1 - \cos^2 \pqty{\frac{\pi}{2}} = 1
    \]
    tells us that \(\sin \pqty{\flatfrac{\pi}{2}} = \pm 1\). But since \(\flatfrac{\pi}{2} \in \oointv{0}{2}\), we have \(\sin \pqty{\flatfrac{\pi}{2}} > 0\), so \(\sin \pqty{\flatfrac{\pi}{2}} = 1\).
\end{proof}

\begin{corollary}[Properties of Complex Exponentials]
    \label{cor:properties-complex-exponentials}%
    \begin{enumerate}
        \item The function \(\exp \pqty{ix}\) for \(x \in \RR\) is periodic, with period \(2 \pi\);
        \item \emph{Euler's Identity}. \(\exp \pqty{i \pi} = -1\);
        \item \(\exp \pqty{i \flatfrac{\pi}{2}} = i\).
    \end{enumerate}
\end{corollary}

Finally, we aim to give \(\pi\) a geometric interpretation.

\begin{lemma}[Perimeter of Unit Circle]
    \label{lem:perimeter-unit-circle}%
    \(2\pi\) is the perimeter of the unit circle
    \[
        \mathcal{S}^1 = \set{z \in \CC}{\abs{z} = 1}.
    \]
\end{lemma}

\begin{proof}
    If we can show the map
    \[
        \function{\gamma}{\cointv{0}{2\pi}}{\mathcal{S}^1}{t}{\exp \pqty{it}}
    \]
    is a bijection, then we get that
    \begin{align*}
        \text{perimeter of }\mathcal{S}^1 &= \text{length of }\gamma \\
        &= \int_{0}^{2 \pi} \abs{\gamma' \pqty{t}} \dd{t} \\
        &= \int_{0}^{2 \pi} \abs{i \exp \pqty{it}} \dd{t} \\
        &= \int_{0}^{2 \pi} \dd{t} \\
        &= 2 \pi.
    \end{align*}

    \textbf{Surjectivity.} Take \(z \in \mathcal{S}^1\), and we write \(z = a + ib\), with \(a^2 + b^2 = 1\), with \(a, b \in \RR\). The restriction
    \[
        \functiondomain{\cos}{\ccintv{0}{\pi}}{\ccintv{-1}{1}}
    \]
    is continuous and strictly monotone. Hence, there is some \(t \in \ccintv{0}{\pi}\) with \(\cos t = a\).

    Furthermore, the restriction
    \[
        \functiondomain{\sin}{\ccintv{0}{\pi}}{\ccintv{0}{1}}
    \]
    is non-negative, so \(\sin t = \sqrt{1 - a^2}\).
    
    \begin{itemize}
        \item If \(b = \sqrt{1 - a^2}\), we get \(\sin t = b\).
        \item If \(b = - \sqrt{1 - a^2}\), we have \(b = - \sin t\), which is to say \(b = \sin \pqty{2 \pi - t}\). But \(\cos t = \cos \pqty{-t} = \cos \pqty{2\pi - t}\).
    \end{itemize}

    Hence, there is some \(t \in \cointv{0}{2\pi}\), such that \(z = \exp \pqty{it}\).

    \textbf{Injectivity.} This is true since \autoref{cor:properties-complex-exponentials} tells us that \(\exp\pqty{it}\) has a least positive period of \(2 \pi\), and so it is injective when restricted to \(\cointv{0}{2\pi}\).
\end{proof}

This allows us now to define the usual trigonometric functions \(\tan, \arccos\) \dots

\subsection{Hyperbolic Trigonometric Functions}

\begin{definition}[Hyperbolic Trigonometric Functions]
    We define the \emph{hyperbolic cosine} as
    \[
        \function{\cosh}{\CC}{\CC}{z}{\frac{1}{2} \bqty{\exp \pqty{z} + \exp \pqty{-z}} = \cos \pqty{iz},}
    \]
    and the \emph{hyperbolic sine} as 
    \[
        \function{\sinh}{\CC}{\CC}{z}{\frac{1}{2} \bqty{\exp \pqty{z} - \exp \pqty{-z}} = - i \sin \pqty{iz}.}
    \]
\end{definition}

We may show that
\begin{itemize}
    \item \(\cosh' z = \sinh z\);
    \item \(\sinh' z = \cosh z\);
    \item \(\cosh^2 z - \sinh^2 z = 1\).
\end{itemize}